%%%%%%%%%%%%%%%%%%%%%%%%%%%%%%%%%%%%%%%%%%%%%%%%%%%%%%%%%%%%
% 10.5 INTEGER PROGRAMMING
% The Composition of Advanced Mathematics
%%%%%%%%%%%%%%%%%%%%%%%%%%%%%%%%%%%%%%%%%%%%%%%%%%%%%%%%%%%%

\section{Integer Programming}
\label{sec:integer-programming}

\prerequisite{
Optimization fundamentals, linear programming, convex optimization,
linear algebra, systems of inequalities, polyhedra, LP duality,
basic feasible solutions, and mathematical modeling.
}

\learningobjectives{
By the end of this section, the reader should be able to:
\begin{itemize}
    \item define integer programming;
    \item distinguish pure integer, binary, and mixed-integer programs;
    \item explain why integrality fundamentally changes optimization;
    \item formulate logical decisions using binary variables;
    \item model either-or conditions and fixed-charge decisions;
    \item formulate assignment, selection, knapsack, facility-location,
          and basic scheduling models;
    \item define the LP relaxation of an integer program;
    \item use LP relaxations to produce bounds;
    \item understand the integrality gap;
    \item recognize integer hulls;
    \item understand why rounding an LP solution may fail;
    \item explain branch-and-bound;
    \item construct branching subproblems;
    \item use incumbent solutions and node bounds;
    \item understand pruning by infeasibility, bounds, and integrality;
    \item distinguish depth-first and best-bound search conceptually;
    \item understand optimality gaps;
    \item recognize cutting planes;
    \item understand valid inequalities;
    \item recognize Gomory cuts conceptually;
    \item formulate cover inequalities for binary knapsack constraints;
    \item understand branch-and-cut;
    \item understand Big-\(M\) formulations and their numerical dangers;
    \item recognize indicator constraints conceptually;
    \item formulate implications using binary variables;
    \item model mutually exclusive choices;
    \item model cardinality constraints;
    \item model piecewise and fixed-cost decisions conceptually;
    \item understand symmetry in integer programs;
    \item recognize symmetry-breaking constraints;
    \item understand preprocessing and bound tightening;
    \item recognize total unimodularity and automatic integrality;
    \item understand Lagrangian relaxation conceptually;
    \item understand decomposition ideas conceptually;
    \item distinguish heuristic solutions from optimality certificates;
    \item understand primal and dual bounds in mixed-integer optimization;
    \item interpret solver optimality gaps;
    \item connect integer programming with combinatorial optimization,
          logistics, scheduling, network design, and operations research.
\end{itemize}
}

%%%%%%%%%%%%%%%%%%%%%%%%%%%%%%%%%%%%%%%%%%%%%%%%%%%%%%%%%%%%
\subsection{What Is Integer Programming?}
\label{subsec:what-is-integer-programming}
%%%%%%%%%%%%%%%%%%%%%%%%%%%%%%%%%%%%%%%%%%%%%%%%%%%%%%%%%%%%

An integer program is an optimization problem in which some or all
decision variables are required to take integer values.

A basic integer linear program has form

\[
\boxed{
\begin{aligned}
\min_{\mathbf{x}}
\quad&
\mathbf{c}^T\mathbf{x}
\\
\text{subject to}
\quad&
A\mathbf{x}\leq\mathbf{b},
\\
&
\mathbf{x}\in\mathbb{Z}^n.
\end{aligned}
}
\]

The objective and constraints may be linear while the integrality
restriction makes the feasible set discrete.

%%%%%%%%%%%%%%%%%%%%%%%%%%%%%%%%%%%%%%%%%%%%%%%%%%%%%%%%%%%%
\subsection{Why Integrality Changes the Problem}
\label{subsec:why-integrality-changes-problem}
%%%%%%%%%%%%%%%%%%%%%%%%%%%%%%%%%%%%%%%%%%%%%%%%%%%%%%%%%%%%

In linear programming, one may move continuously through a convex
polyhedron.

In integer programming, only selected lattice points are feasible.

Thus

\[
\boxed{
\text{continuous feasible region}
}
\]

is replaced by

\[
\boxed{
\text{discrete feasible points}.
}
\]

The resulting feasible set is generally nonconvex.

%%%%%%%%%%%%%%%%%%%%%%%%%%%%%%%%%%%%%%%%%%%%%%%%%%%%%%%%%%%%
\subsection{Pure Integer Programming}
\label{subsec:pure-integer-programming}
%%%%%%%%%%%%%%%%%%%%%%%%%%%%%%%%%%%%%%%%%%%%%%%%%%%%%%%%%%%%

A pure integer program requires every decision variable to be integral:

\[
\boxed{
x_j\in\mathbb{Z},
\qquad
j=1,\ldots,n.
}
\]

For example,

\[
\boxed{
\begin{aligned}
\max
\quad&
4x_1+7x_2
\\
\text{subject to}
\quad&
3x_1+2x_2\leq20,
\\
&
x_1,x_2\in\mathbb{Z}_{\geq0}.
\end{aligned}
}
\]

%%%%%%%%%%%%%%%%%%%%%%%%%%%%%%%%%%%%%%%%%%%%%%%%%%%%%%%%%%%%
\subsection{Binary Integer Programming}
\label{subsec:binary-integer-programming}
%%%%%%%%%%%%%%%%%%%%%%%%%%%%%%%%%%%%%%%%%%%%%%%%%%%%%%%%%%%%

A binary variable satisfies

\[
\boxed{
x_j\in\{0,1\}.
}
\]

Binary variables naturally encode:

\[
\boxed{
0=\text{no},
\qquad
1=\text{yes}.
}
\]

They are used for:

\[
\boxed{
\text{selection},
}
\]

\[
\boxed{
\text{activation},
}
\]

\[
\boxed{
\text{assignment},
}
\]

and

\[
\boxed{
\text{logical decisions}.
}
\]

%%%%%%%%%%%%%%%%%%%%%%%%%%%%%%%%%%%%%%%%%%%%%%%%%%%%%%%%%%%%
\subsection{Mixed-Integer Programming}
\label{subsec:mixed-integer-programming}
%%%%%%%%%%%%%%%%%%%%%%%%%%%%%%%%%%%%%%%%%%%%%%%%%%%%%%%%%%%%

A mixed-integer program contains both integer and continuous variables.

A mixed-integer linear program, abbreviated MILP, may be written

\[
\boxed{
\begin{aligned}
\min_{\mathbf{x},\mathbf{y}}
\quad&
\mathbf{c}^T\mathbf{x}
+
\mathbf{d}^T\mathbf{y}
\\
\text{subject to}
\quad&
A\mathbf{x}
+
B\mathbf{y}
\leq
\mathbf{b},
\\
&
\mathbf{x}\in\mathbb{Z}^p,
\\
&
\mathbf{y}\in\mathbb{R}^q.
\end{aligned}
}
\]

MILPs are among the most widely used operations-research models.

%%%%%%%%%%%%%%%%%%%%%%%%%%%%%%%%%%%%%%%%%%%%%%%%%%%%%%%%%%%%
\subsection{Integer Lattice}
\label{subsec:integer-lattice}
%%%%%%%%%%%%%%%%%%%%%%%%%%%%%%%%%%%%%%%%%%%%%%%%%%%%%%%%%%%%

The integer lattice in \(\mathbb{R}^n\) is

\[
\boxed{
\mathbb{Z}^n.
}
\]

An integer-programming feasible set can often be viewed as

\[
\boxed{
P\cap\mathbb{Z}^n,
}
\]

where \(P\) is a polyhedron.

The continuous polyhedron provides geometric information, while the
integer lattice imposes discreteness.

%%%%%%%%%%%%%%%%%%%%%%%%%%%%%%%%%%%%%%%%%%%%%%%%%%%%%%%%%%%%
\subsection{LP Relaxation}
\label{subsec:integer-lp-relaxation}
%%%%%%%%%%%%%%%%%%%%%%%%%%%%%%%%%%%%%%%%%%%%%%%%%%%%%%%%%%%%

The LP relaxation of an integer program is obtained by removing the
integrality restrictions.

For

\[
\boxed{
\begin{aligned}
\max
\quad&
c^Tx
\\
\text{subject to}
\quad&
Ax\leq b,
\\
&
x\in\mathbb{Z}_{\geq0}^n,
\end{aligned}
}
\]

the LP relaxation is

\[
\boxed{
\begin{aligned}
\max
\quad&
c^Tx
\\
\text{subject to}
\quad&
Ax\leq b,
\\
&
x\geq0.
\end{aligned}
}
\]

%%%%%%%%%%%%%%%%%%%%%%%%%%%%%%%%%%%%%%%%%%%%%%%%%%%%%%%%%%%%
\subsection{Relaxation Bounds}
\label{subsec:integer-relaxation-bounds}
%%%%%%%%%%%%%%%%%%%%%%%%%%%%%%%%%%%%%%%%%%%%%%%%%%%%%%%%%%%%

For a maximization problem, the LP relaxation has a larger feasible set.

Therefore,

\[
\boxed{
z_{\mathrm{IP}}^*
\leq
z_{\mathrm{LP}}^*.
}
\]

Thus the LP relaxation provides an upper bound.

For a minimization problem,

\[
\boxed{
z_{\mathrm{LP}}^*
\leq
z_{\mathrm{IP}}^*.
}
\]

Thus the LP relaxation provides a lower bound.

%%%%%%%%%%%%%%%%%%%%%%%%%%%%%%%%%%%%%%%%%%%%%%%%%%%%%%%%%%%%
\subsection{Worked Example: LP Relaxation}
\label{subsec:worked-ip-relaxation}
%%%%%%%%%%%%%%%%%%%%%%%%%%%%%%%%%%%%%%%%%%%%%%%%%%%%%%%%%%%%

\begin{workedexamplebox}{Fractional LP Solution}

Consider

\[
\begin{aligned}
\max
\quad&
x+y
\\
\text{subject to}
\quad&
2x+2y\leq5,
\\
&
x,y\in\mathbb{Z}_{\geq0}.
\end{aligned}
\]

The LP relaxation allows real variables:

\[
x,y\geq0.
\]

Its boundary is

\[
x+y\leq\frac52.
\]

Therefore,

\[
z_{\mathrm{LP}}^*
=
\frac52.
\]

For integer variables,

\[
x+y
\]

must be an integer and cannot exceed \(2.5\), so

\[
z_{\mathrm{IP}}^*=2.
\]

Hence,

\begin{answerbox}
\[
\boxed{
z_{\mathrm{LP}}^*
=
2.5,
\qquad
z_{\mathrm{IP}}^*
=
2.
}
\]
\end{answerbox}

\end{workedexamplebox}

%%%%%%%%%%%%%%%%%%%%%%%%%%%%%%%%%%%%%%%%%%%%%%%%%%%%%%%%%%%%
\subsection{Integrality Gap}
\label{subsec:integrality-gap}
%%%%%%%%%%%%%%%%%%%%%%%%%%%%%%%%%%%%%%%%%%%%%%%%%%%%%%%%%%%%

The difference between an integer optimum and the optimum of a relaxation
measures the quality of the relaxation.

For a minimization problem, one may examine

\[
\boxed{
z_{\mathrm{IP}}^*
-
z_{\mathrm{LP}}^*.
}
\]

A relative gap can also be used.

The exact definition depends on whether the problem is minimization or
maximization and on the normalization convention.

%%%%%%%%%%%%%%%%%%%%%%%%%%%%%%%%%%%%%%%%%%%%%%%%%%%%%%%%%%%%
\subsection{Why Strong Relaxations Matter}
\label{subsec:strong-relaxations}
%%%%%%%%%%%%%%%%%%%%%%%%%%%%%%%%%%%%%%%%%%%%%%%%%%%%%%%%%%%%

A relaxation is strong when its optimal value is close to the integer
optimal value.

Strong relaxations help because they:

\[
\boxed{
\text{produce tighter bounds},
}
\]

\[
\boxed{
\text{reduce branch-and-bound search},
}
\]

and

\[
\boxed{
\text{improve optimality certificates}.
}
\]

Formulation quality can therefore be as important as solver choice.

%%%%%%%%%%%%%%%%%%%%%%%%%%%%%%%%%%%%%%%%%%%%%%%%%%%%%%%%%%%%
\subsection{Integer Hull}
\label{subsec:integer-hull}
%%%%%%%%%%%%%%%%%%%%%%%%%%%%%%%%%%%%%%%%%%%%%%%%%%%%%%%%%%%%

Let

\[
S
=
P\cap\mathbb{Z}^n.
\]

The integer hull is

\[
\boxed{
\operatorname{conv}(S).
}
\]

It is the convex hull of all integer feasible points.

Optimizing a linear objective over the integer hull gives the same optimal
value as optimizing over the original integer points.

%%%%%%%%%%%%%%%%%%%%%%%%%%%%%%%%%%%%%%%%%%%%%%%%%%%%%%%%%%%%
\subsection{Ideal Formulations}
\label{subsec:ideal-integer-formulations}
%%%%%%%%%%%%%%%%%%%%%%%%%%%%%%%%%%%%%%%%%%%%%%%%%%%%%%%%%%%%

A formulation is especially strong if the extreme points of its LP
relaxation are already integral.

Such a formulation can be called ideal in an appropriate modeling sense.

Then solving the LP relaxation may already solve the integer problem.

%%%%%%%%%%%%%%%%%%%%%%%%%%%%%%%%%%%%%%%%%%%%%%%%%%%%%%%%%%%%
\subsection{Why Simple Rounding Can Fail}
\label{subsec:why-rounding-can-fail}
%%%%%%%%%%%%%%%%%%%%%%%%%%%%%%%%%%%%%%%%%%%%%%%%%%%%%%%%%%%%

Suppose an LP relaxation gives

\[
x_1=1.7,
\qquad
x_2=2.6.
\]

Rounding gives

\[
x_1=2,
\qquad
x_2=3.
\]

However, the rounded point may violate the constraints.

Even if it remains feasible, it may be far from optimal.

Therefore,

\[
\boxed{
\text{solve LP and round}
}
\]

is not a general integer-programming algorithm.

%%%%%%%%%%%%%%%%%%%%%%%%%%%%%%%%%%%%%%%%%%%%%%%%%%%%%%%%%%%%
\subsection{Worked Example: Rounding Failure}
\label{subsec:worked-rounding-failure}
%%%%%%%%%%%%%%%%%%%%%%%%%%%%%%%%%%%%%%%%%%%%%%%%%%%%%%%%%%%%

\begin{workedexamplebox}{Rounding Can Destroy Feasibility}

Suppose

\[
x+y\leq3,
\]

with

\[
x,y\in\mathbb{Z}_{\geq0}.
\]

Assume an LP solution is

\[
x=1.6,
\qquad
y=1.4.
\]

This point satisfies

\[
x+y=3.
\]

Rounding each component upward gives

\[
x=2,
\qquad
y=2.
\]

Then

\[
2+2=4>3.
\]

Therefore,

\begin{answerbox}
\[
\boxed{
\text{componentwise rounding can produce an infeasible solution}.
}
\]
\end{answerbox}

\end{workedexamplebox}

%%%%%%%%%%%%%%%%%%%%%%%%%%%%%%%%%%%%%%%%%%%%%%%%%%%%%%%%%%%%
\subsection{Binary Selection Models}
\label{subsec:binary-selection-models}
%%%%%%%%%%%%%%%%%%%%%%%%%%%%%%%%%%%%%%%%%%%%%%%%%%%%%%%%%%%%

Suppose project \(j\) is either selected or not selected.

Define

\[
\boxed{
x_j
=
\begin{cases}
1,
&
\text{project }j\text{ is selected},
\\
0,
&
\text{otherwise}.
\end{cases}
}
\]

If project \(j\) provides value \(v_j\), then total value is

\[
\boxed{
\sum_j
v_jx_j.
}
\]

%%%%%%%%%%%%%%%%%%%%%%%%%%%%%%%%%%%%%%%%%%%%%%%%%%%%%%%%%%%%
\subsection{Cardinality Constraints}
\label{subsec:cardinality-constraints}
%%%%%%%%%%%%%%%%%%%%%%%%%%%%%%%%%%%%%%%%%%%%%%%%%%%%%%%%%%%%

If at most \(k\) items may be selected,

\[
\boxed{
\sum_{j=1}^{n}
x_j
\leq
k,
}
\]

where

\[
x_j\in\{0,1\}.
\]

If exactly \(k\) items must be selected,

\[
\boxed{
\sum_{j=1}^{n}
x_j
=
k.
}
\]

%%%%%%%%%%%%%%%%%%%%%%%%%%%%%%%%%%%%%%%%%%%%%%%%%%%%%%%%%%%%
\subsection{At Least and At Most One}
\label{subsec:at-least-at-most-one}
%%%%%%%%%%%%%%%%%%%%%%%%%%%%%%%%%%%%%%%%%%%%%%%%%%%%%%%%%%%%

For binary variables:

At most one choice:

\[
\boxed{
\sum_jx_j\leq1.
}
\]

Exactly one choice:

\[
\boxed{
\sum_jx_j=1.
}
\]

At least one choice:

\[
\boxed{
\sum_jx_j\geq1.
}
\]

These simple constraints encode many logical decisions.

%%%%%%%%%%%%%%%%%%%%%%%%%%%%%%%%%%%%%%%%%%%%%%%%%%%%%%%%%%%%
\subsection{Logical Implication}
\label{subsec:binary-logical-implication}
%%%%%%%%%%%%%%%%%%%%%%%%%%%%%%%%%%%%%%%%%%%%%%%%%%%%%%%%%%%%

Suppose

\[
x,y\in\{0,1\},
\]

and selecting \(x\) requires selecting \(y\).

The implication

\[
\boxed{
x=1
\Rightarrow
y=1
}
\]

is modeled by

\[
\boxed{
x\leq y.
}
\]

%%%%%%%%%%%%%%%%%%%%%%%%%%%%%%%%%%%%%%%%%%%%%%%%%%%%%%%%%%%%
\subsection{Mutual Exclusion}
\label{subsec:binary-mutual-exclusion}
%%%%%%%%%%%%%%%%%%%%%%%%%%%%%%%%%%%%%%%%%%%%%%%%%%%%%%%%%%%%

Suppose choices \(x\) and \(y\) cannot both occur.

Then

\[
\boxed{
x+y\leq1.
}
\]

For several mutually exclusive options,

\[
\boxed{
\sum_jx_j\leq1.
}
\]

%%%%%%%%%%%%%%%%%%%%%%%%%%%%%%%%%%%%%%%%%%%%%%%%%%%%%%%%%%%%
\subsection{Either-Or Logic}
\label{subsec:either-or-logic}
%%%%%%%%%%%%%%%%%%%%%%%%%%%%%%%%%%%%%%%%%%%%%%%%%%%%%%%%%%%%

Suppose exactly one of two alternatives must be chosen.

Using

\[
y\in\{0,1\},
\]

one may let:

\[
\boxed{
y=0
\rightarrow
\text{alternative A},
}
\]

and

\[
\boxed{
y=1
\rightarrow
\text{alternative B}.
}
\]

Conditional constraints can then be activated using indicator constraints
or sufficiently strong Big-\(M\) formulations.

%%%%%%%%%%%%%%%%%%%%%%%%%%%%%%%%%%%%%%%%%%%%%%%%%%%%%%%%%%%%
\subsection{Big-\(M\) Implication}
\label{subsec:big-m-implication}
%%%%%%%%%%%%%%%%%%%%%%%%%%%%%%%%%%%%%%%%%%%%%%%%%%%%%%%%%%%%

Suppose

\[
y\in\{0,1\}
\]

and we want

\[
\boxed{
y=1
\Rightarrow
a^Tx\leq b.
}
\]

A Big-\(M\) formulation is

\[
\boxed{
a^Tx
\leq
b
+
M(1-y).
}
\]

If

\[
y=1,
\]

the original constraint is enforced.

If

\[
y=0,
\]

the constraint becomes relaxed by \(M\).

%%%%%%%%%%%%%%%%%%%%%%%%%%%%%%%%%%%%%%%%%%%%%%%%%%%%%%%%%%%%
\subsection{Choosing Big-\(M\)}
\label{subsec:choosing-big-m}
%%%%%%%%%%%%%%%%%%%%%%%%%%%%%%%%%%%%%%%%%%%%%%%%%%%%%%%%%%%%

The constant \(M\) must be large enough to make the constraint inactive
when intended.

However, choosing \(M\) excessively large can cause:

\[
\boxed{
\text{weak LP relaxations},
}
\]

\[
\boxed{
\text{poor numerical conditioning},
}
\]

and

\[
\boxed{
\text{slow branch-and-bound}.
}
\]

Therefore one should use the smallest valid \(M\) that can be justified.

%%%%%%%%%%%%%%%%%%%%%%%%%%%%%%%%%%%%%%%%%%%%%%%%%%%%%%%%%%%%
\subsection{Indicator Constraints}
\label{subsec:indicator-constraints}
%%%%%%%%%%%%%%%%%%%%%%%%%%%%%%%%%%%%%%%%%%%%%%%%%%%%%%%%%%%%

Many modern solvers support logical conditions directly:

\[
\boxed{
y=1
\Rightarrow
a^Tx\leq b.
}
\]

These are called indicator constraints.

They can avoid manually selecting Big-\(M\) values, although the solver
still needs an internal mathematical representation.

%%%%%%%%%%%%%%%%%%%%%%%%%%%%%%%%%%%%%%%%%%%%%%%%%%%%%%%%%%%%
\subsection{Fixed-Charge Models}
\label{subsec:fixed-charge-models}
%%%%%%%%%%%%%%%%%%%%%%%%%%%%%%%%%%%%%%%%%%%%%%%%%%%%%%%%%%%%

Suppose using an activity incurs:

\[
\boxed{
\text{fixed cost }F
}
\]

plus

\[
\boxed{
\text{variable cost }cx.
}
\]

Let

\[
y\in\{0,1\}
\]

indicate whether the activity is used.

Then the objective contains

\[
\boxed{
Fy+cx.
}
\]

To link \(x\) and \(y\), use

\[
\boxed{
0\leq x\leq Uy,
}
\]

where \(U\) is a valid upper bound.

%%%%%%%%%%%%%%%%%%%%%%%%%%%%%%%%%%%%%%%%%%%%%%%%%%%%%%%%%%%%
\subsection{Worked Example: Fixed Setup Cost}
\label{subsec:worked-fixed-charge}
%%%%%%%%%%%%%%%%%%%%%%%%%%%%%%%%%%%%%%%%%%%%%%%%%%%%%%%%%%%%

\begin{workedexamplebox}{Activating Production}

Suppose production quantity \(x\) has variable cost \(3x\) and a setup
cost of \(100\) if any production occurs.

Let

\[
y\in\{0,1\}
\]

indicate whether production is active.

If capacity is at most \(50\), impose

\[
0\leq x\leq50y.
\]

The total cost is

\[
3x+100y.
\]

Thus,

\begin{answerbox}
\[
\boxed{
\begin{aligned}
\min_{x,y}
\quad&
3x+100y
\\
\text{subject to}
\quad&
0\leq x\leq50y,
\\
&
y\in\{0,1\}.
\end{aligned}
}
\]
\end{answerbox}

If \(x>0\), then necessarily \(y=1\).

\end{workedexamplebox}

%%%%%%%%%%%%%%%%%%%%%%%%%%%%%%%%%%%%%%%%%%%%%%%%%%%%%%%%%%%%
\subsection{Knapsack Problem}
\label{subsec:integer-knapsack-problem}
%%%%%%%%%%%%%%%%%%%%%%%%%%%%%%%%%%%%%%%%%%%%%%%%%%%%%%%%%%%%

Suppose item \(j\) has:

\[
v_j
=
\text{value},
\]

and

\[
w_j
=
\text{weight}.
\]

Let

\[
x_j\in\{0,1\}
\]

indicate whether the item is selected.

Then the binary knapsack problem is

\[
\boxed{
\begin{aligned}
\max
\quad&
\sum_{j=1}^{n}
v_jx_j
\\
\text{subject to}
\quad&
\sum_{j=1}^{n}
w_jx_j
\leq
W,
\\
&
x_j\in\{0,1\}.
\end{aligned}
}
\]

%%%%%%%%%%%%%%%%%%%%%%%%%%%%%%%%%%%%%%%%%%%%%%%%%%%%%%%%%%%%
\subsection{Worked Example: Binary Knapsack}
\label{subsec:worked-binary-knapsack}
%%%%%%%%%%%%%%%%%%%%%%%%%%%%%%%%%%%%%%%%%%%%%%%%%%%%%%%%%%%%

\begin{workedexamplebox}{Selecting Projects}

Suppose three projects have values

\[
8,\ 11,\ 6
\]

and costs

\[
4,\ 6,\ 3.
\]

The budget is \(9\).

Let

\[
x_j\in\{0,1\}.
\]

Then

\[
\boxed{
\begin{aligned}
\max
\quad&
8x_1+11x_2+6x_3
\\
\text{subject to}
\quad&
4x_1+6x_2+3x_3
\leq9.
\end{aligned}
}
\]

Checking feasible combinations:

\[
(1,0,1)
\]

has cost \(7\) and value \(14\).

\[
(0,1,1)
\]

has cost \(9\) and value \(17\).

\[
(1,1,0)
\]

has cost \(10\) and is infeasible.

Therefore,

\begin{answerbox}
\[
\boxed{
(x_1,x_2,x_3)
=
(0,1,1)
}
\]

with

\[
\boxed{
z^*=17.
}
\]
\end{answerbox}

\end{workedexamplebox}

%%%%%%%%%%%%%%%%%%%%%%%%%%%%%%%%%%%%%%%%%%%%%%%%%%%%%%%%%%%%
\subsection{Assignment Problem}
\label{subsec:integer-assignment-problem}
%%%%%%%%%%%%%%%%%%%%%%%%%%%%%%%%%%%%%%%%%%%%%%%%%%%%%%%%%%%%

Suppose workers \(i\) must be assigned to jobs \(j\).

Define

\[
\boxed{
x_{ij}
=
\begin{cases}
1,
&
\text{worker }i\text{ performs job }j,
\\
0,
&
\text{otherwise}.
\end{cases}
}
\]

If \(c_{ij}\) is assignment cost, then

\[
\boxed{
\min
\sum_i
\sum_j
c_{ij}x_{ij}.
}
\]

%%%%%%%%%%%%%%%%%%%%%%%%%%%%%%%%%%%%%%%%%%%%%%%%%%%%%%%%%%%%
\subsection{Assignment Constraints}
\label{subsec:assignment-constraints}
%%%%%%%%%%%%%%%%%%%%%%%%%%%%%%%%%%%%%%%%%%%%%%%%%%%%%%%%%%%%

If each worker receives exactly one job,

\[
\boxed{
\sum_j
x_{ij}
=
1
}
\]

for every worker \(i\).

If each job receives exactly one worker,

\[
\boxed{
\sum_i
x_{ij}
=
1
}
\]

for every job \(j\).

Thus,

\[
\boxed{
x_{ij}\in\{0,1\}.
}
\]

The assignment problem has special structure that often makes its LP
relaxation integral.

%%%%%%%%%%%%%%%%%%%%%%%%%%%%%%%%%%%%%%%%%%%%%%%%%%%%%%%%%%%%
\subsection{Facility Location}
\label{subsec:facility-location}
%%%%%%%%%%%%%%%%%%%%%%%%%%%%%%%%%%%%%%%%%%%%%%%%%%%%%%%%%%%%

Suppose potential facility \(j\) has fixed opening cost \(F_j\).

Define

\[
\boxed{
y_j
=
\begin{cases}
1,
&
\text{facility }j\text{ is opened},
\\
0,
&
\text{otherwise}.
\end{cases}
}
\]

Let

\[
x_{ij}
\]

represent service from facility \(j\) to customer \(i\).

Linking constraints ensure customers can only be served from open
facilities.

A common form is

\[
\boxed{
x_{ij}\leq y_j.
}
\]

%%%%%%%%%%%%%%%%%%%%%%%%%%%%%%%%%%%%%%%%%%%%%%%%%%%%%%%%%%%%
\subsection{Capacitated Facility Location}
\label{subsec:capacitated-facility-location}
%%%%%%%%%%%%%%%%%%%%%%%%%%%%%%%%%%%%%%%%%%%%%%%%%%%%%%%%%%%%

If facility \(j\) has capacity \(U_j\),

\[
\boxed{
\sum_i
d_i x_{ij}
\leq
U_j y_j.
}
\]

When

\[
y_j=0,
\]

the right side is zero and no demand can be assigned there.

When

\[
y_j=1,
\]

capacity becomes available.

%%%%%%%%%%%%%%%%%%%%%%%%%%%%%%%%%%%%%%%%%%%%%%%%%%%%%%%%%%%%
\subsection{Basic Scheduling Variables}
\label{subsec:integer-scheduling-variables}
%%%%%%%%%%%%%%%%%%%%%%%%%%%%%%%%%%%%%%%%%%%%%%%%%%%%%%%%%%%%

Binary variables can encode whether job \(j\) is assigned to machine
\(m\):

\[
\boxed{
x_{jm}
\in
\{0,1\}.
}
\]

An assignment constraint may be

\[
\boxed{
\sum_m
x_{jm}
=
1.
}
\]

Scheduling models may also require integer start times, precedence
relations, or sequencing variables.

More detailed scheduling theory appears later in Chapter 10.

%%%%%%%%%%%%%%%%%%%%%%%%%%%%%%%%%%%%%%%%%%%%%%%%%%%%%%%%%%%%
\subsection{Precedence Constraints}
\label{subsec:integer-precedence-constraints}
%%%%%%%%%%%%%%%%%%%%%%%%%%%%%%%%%%%%%%%%%%%%%%%%%%%%%%%%%%%%

If job \(i\) must finish before job \(j\) starts, and

\[
s_i
=
\text{start time of job }i,
\]

with processing time \(p_i\), then

\[
\boxed{
s_j
\geq
s_i+p_i.
}
\]

If \(s_i\) are continuous, the model can still be mixed-integer because
binary variables may determine sequencing decisions.

%%%%%%%%%%%%%%%%%%%%%%%%%%%%%%%%%%%%%%%%%%%%%%%%%%%%%%%%%%%%
\subsection{Sequencing with Binary Variables}
\label{subsec:sequencing-binary-variables}
%%%%%%%%%%%%%%%%%%%%%%%%%%%%%%%%%%%%%%%%%%%%%%%%%%%%%%%%%%%%

Suppose jobs \(i\) and \(j\) cannot overlap.

Define

\[
y_{ij}
=
1
\]

if \(i\) occurs before \(j\).

A Big-\(M\) formulation can use

\[
\boxed{
s_i+p_i
\leq
s_j+M(1-y_{ij}),
}
\]

and

\[
\boxed{
s_j+p_j
\leq
s_i+My_{ij}.
}
\]

Exactly one ordering is enforced when \(M\) is valid.

%%%%%%%%%%%%%%%%%%%%%%%%%%%%%%%%%%%%%%%%%%%%%%%%%%%%%%%%%%%%
\subsection{Branch-and-Bound}
\label{subsec:branch-and-bound}
%%%%%%%%%%%%%%%%%%%%%%%%%%%%%%%%%%%%%%%%%%%%%%%%%%%%%%%%%%%%

Branch-and-bound is a general method for solving integer programs.

It combines:

\[
\boxed{
\text{relaxation},
}
\]

\[
\boxed{
\text{branching},
}
\]

and

\[
\boxed{
\text{bounding}.
}
\]

The method systematically partitions the feasible set while using bounds
to eliminate regions that cannot contain an improved solution.

%%%%%%%%%%%%%%%%%%%%%%%%%%%%%%%%%%%%%%%%%%%%%%%%%%%%%%%%%%%%
\subsection{Root Relaxation}
\label{subsec:branch-bound-root}
%%%%%%%%%%%%%%%%%%%%%%%%%%%%%%%%%%%%%%%%%%%%%%%%%%%%%%%%%%%%

The process usually begins by solving the LP relaxation.

This is the root node of the search tree.

If the LP solution is already integral, then the integer problem may be
solved immediately.

Otherwise, one selects a fractional integer variable for branching.

%%%%%%%%%%%%%%%%%%%%%%%%%%%%%%%%%%%%%%%%%%%%%%%%%%%%%%%%%%%%
\subsection{Branching on a Fractional Variable}
\label{subsec:branching-fractional-variable}
%%%%%%%%%%%%%%%%%%%%%%%%%%%%%%%%%%%%%%%%%%%%%%%%%%%%%%%%%%%%

Suppose the LP relaxation gives

\[
\boxed{
x_j^*=3.7,
}
\]

where \(x_j\) must be integer.

Then every integer solution satisfies either

\[
\boxed{
x_j\leq3
}
\]

or

\[
\boxed{
x_j\geq4.
}
\]

Thus create two child nodes:

\[
\boxed{
x_j\leq\lfloor3.7\rfloor=3,
}
\]

and

\[
\boxed{
x_j\geq\lceil3.7\rceil=4.
}
\]

%%%%%%%%%%%%%%%%%%%%%%%%%%%%%%%%%%%%%%%%%%%%%%%%%%%%%%%%%%%%
\subsection{Branch-and-Bound Tree}
\label{subsec:branch-and-bound-tree}
%%%%%%%%%%%%%%%%%%%%%%%%%%%%%%%%%%%%%%%%%%%%%%%%%%%%%%%%%%%%

Each node represents the original model plus additional branching
constraints.

Solving the node relaxation gives:

\[
\boxed{
\text{a bound},
}
\]

and sometimes

\[
\boxed{
\text{an integer feasible solution}.
}
\]

The tree continues only where further exploration may improve the current
best integer solution.

%%%%%%%%%%%%%%%%%%%%%%%%%%%%%%%%%%%%%%%%%%%%%%%%%%%%%%%%%%%%
\subsection{Incumbent Solution}
\label{subsec:integer-incumbent}
%%%%%%%%%%%%%%%%%%%%%%%%%%%%%%%%%%%%%%%%%%%%%%%%%%%%%%%%%%%%

The best integer feasible solution found so far is called the incumbent.

For minimization, let

\[
\boxed{
U
=
\text{incumbent objective value}.
}
\]

Then \(U\) is an upper bound on the unknown optimal integer objective.

For maximization, the incumbent provides a lower bound.

%%%%%%%%%%%%%%%%%%%%%%%%%%%%%%%%%%%%%%%%%%%%%%%%%%%%%%%%%%%%
\subsection{Node Bounds}
\label{subsec:integer-node-bounds}
%%%%%%%%%%%%%%%%%%%%%%%%%%%%%%%%%%%%%%%%%%%%%%%%%%%%%%%%%%%%

For a minimization problem, the LP relaxation at a node gives a lower
bound:

\[
\boxed{
L_{\text{node}}
\leq
z_{\text{IP,node}}^*.
}
\]

If

\[
L_{\text{node}}
\geq
U,
\]

where \(U\) is the incumbent value, that node cannot improve the incumbent.

It may therefore be pruned.

%%%%%%%%%%%%%%%%%%%%%%%%%%%%%%%%%%%%%%%%%%%%%%%%%%%%%%%%%%%%
\subsection{Pruning by Infeasibility}
\label{subsec:pruning-infeasibility}
%%%%%%%%%%%%%%%%%%%%%%%%%%%%%%%%%%%%%%%%%%%%%%%%%%%%%%%%%%%%

If a node's LP relaxation is infeasible, then the corresponding integer
subproblem is also infeasible.

Therefore the node can be removed immediately.

%%%%%%%%%%%%%%%%%%%%%%%%%%%%%%%%%%%%%%%%%%%%%%%%%%%%%%%%%%%%
\subsection{Pruning by Integrality}
\label{subsec:pruning-integrality}
%%%%%%%%%%%%%%%%%%%%%%%%%%%%%%%%%%%%%%%%%%%%%%%%%%%%%%%%%%%%

If the LP relaxation solution at a node satisfies all required integrality
conditions, then it is feasible for the integer problem.

The incumbent can be updated if the objective is better.

No further branching is needed at that node.

%%%%%%%%%%%%%%%%%%%%%%%%%%%%%%%%%%%%%%%%%%%%%%%%%%%%%%%%%%%%
\subsection{Pruning by Bound}
\label{subsec:pruning-bound}
%%%%%%%%%%%%%%%%%%%%%%%%%%%%%%%%%%%%%%%%%%%%%%%%%%%%%%%%%%%%

For minimization, suppose the best known integer objective is \(U\).

If a node relaxation gives

\[
\boxed{
L_{\text{node}}\geq U,
}
\]

then every integer solution in that node has objective at least \(U\).

The node cannot improve the incumbent and is pruned.

%%%%%%%%%%%%%%%%%%%%%%%%%%%%%%%%%%%%%%%%%%%%%%%%%%%%%%%%%%%%
\subsection{Worked Example: Branching}
\label{subsec:worked-branching}
%%%%%%%%%%%%%%%%%%%%%%%%%%%%%%%%%%%%%%%%%%%%%%%%%%%%%%%%%%%%

\begin{workedexamplebox}{Creating Two Child Nodes}

Suppose an LP relaxation of an integer minimization problem gives

\[
x_1=2.4.
\]

Because

\[
x_1\in\mathbb{Z},
\]

branch on \(x_1\).

The first child adds

\[
x_1\leq2.
\]

The second child adds

\[
x_1\geq3.
\]

Thus,

\begin{answerbox}
\[
\boxed{
\text{Child 1: }x_1\leq2,
}
\]

\[
\boxed{
\text{Child 2: }x_1\geq3.
}
\]
\end{answerbox}

Every integer feasible point belongs to one of these two branches.

\end{workedexamplebox}

%%%%%%%%%%%%%%%%%%%%%%%%%%%%%%%%%%%%%%%%%%%%%%%%%%%%%%%%%%%%
\subsection{Global Bounds}
\label{subsec:integer-global-bounds}
%%%%%%%%%%%%%%%%%%%%%%%%%%%%%%%%%%%%%%%%%%%%%%%%%%%%%%%%%%%%

For minimization, suppose:

\[
U
=
\text{best incumbent objective},
\]

and

\[
L
=
\text{best valid lower bound from unexplored nodes}.
\]

Then

\[
\boxed{
L
\leq
z_{\mathrm{IP}}^*
\leq
U.
}
\]

As the search progresses, the interval narrows.

%%%%%%%%%%%%%%%%%%%%%%%%%%%%%%%%%%%%%%%%%%%%%%%%%%%%%%%%%%%%
\subsection{Absolute Optimality Gap}
\label{subsec:absolute-mip-gap}
%%%%%%%%%%%%%%%%%%%%%%%%%%%%%%%%%%%%%%%%%%%%%%%%%%%%%%%%%%%%

For minimization, an absolute gap is

\[
\boxed{
U-L.
}
\]

If

\[
\boxed{
U-L=0,
}
\]

then the incumbent is proven globally optimal.

Numerical solvers generally use tolerances rather than requiring exact
floating-point equality.

%%%%%%%%%%%%%%%%%%%%%%%%%%%%%%%%%%%%%%%%%%%%%%%%%%%%%%%%%%%%
\subsection{Relative Optimality Gap}
\label{subsec:relative-mip-gap}
%%%%%%%%%%%%%%%%%%%%%%%%%%%%%%%%%%%%%%%%%%%%%%%%%%%%%%%%%%%%

A common relative-gap form is

\[
\boxed{
\frac{
|U-L|
}{
\max
\{
1,|U|
\}
}.
}
\]

Exact solver conventions vary.

A solver may stop when the relative gap is below a user-specified
tolerance.

%%%%%%%%%%%%%%%%%%%%%%%%%%%%%%%%%%%%%%%%%%%%%%%%%%%%%%%%%%%%
\subsection{Node Selection}
\label{subsec:branch-bound-node-selection}
%%%%%%%%%%%%%%%%%%%%%%%%%%%%%%%%%%%%%%%%%%%%%%%%%%%%%%%%%%%%

A branch-and-bound algorithm must decide which open node to explore next.

Common strategies include:

\[
\boxed{
\text{depth-first search},
}
\]

\[
\boxed{
\text{breadth-oriented strategies},
}
\]

and

\[
\boxed{
\text{best-bound search}.
}
\]

Different strategies trade memory use, incumbent discovery, and bound
improvement.

%%%%%%%%%%%%%%%%%%%%%%%%%%%%%%%%%%%%%%%%%%%%%%%%%%%%%%%%%%%%
\subsection{Branching Rules}
\label{subsec:branching-rules}
%%%%%%%%%%%%%%%%%%%%%%%%%%%%%%%%%%%%%%%%%%%%%%%%%%%%%%%%%%%%

The choice of branching variable can strongly affect search-tree size.

Possible strategies include:

\[
\boxed{
\text{most fractional branching},
}
\]

\[
\boxed{
\text{strong branching},
}
\]

and

\[
\boxed{
\text{pseudocost branching}.
}
\]

Modern solvers combine several sophisticated branching ideas.

%%%%%%%%%%%%%%%%%%%%%%%%%%%%%%%%%%%%%%%%%%%%%%%%%%%%%%%%%%%%
\subsection{Strong Branching Preview}
\label{subsec:strong-branching-preview}
%%%%%%%%%%%%%%%%%%%%%%%%%%%%%%%%%%%%%%%%%%%%%%%%%%%%%%%%%%%%

Strong branching temporarily evaluates several possible branch choices by
partially solving their child relaxations.

The variable producing the most promising bound improvement is selected.

This can reduce the search tree but increases work per branching decision.

%%%%%%%%%%%%%%%%%%%%%%%%%%%%%%%%%%%%%%%%%%%%%%%%%%%%%%%%%%%%
\subsection{Cutting Planes}
\label{subsec:cutting-planes}
%%%%%%%%%%%%%%%%%%%%%%%%%%%%%%%%%%%%%%%%%%%%%%%%%%%%%%%%%%%%

A cutting plane is a valid inequality that removes fractional solutions
without removing any feasible integer solution.

Suppose the current LP relaxation contains a fractional optimum.

Add an inequality

\[
\boxed{
a^Tx\leq b
}
\]

such that:

\[
\boxed{
\text{every integer feasible point satisfies it},
}
\]

but

\[
\boxed{
\text{the current fractional point violates it}.
}
\]

%%%%%%%%%%%%%%%%%%%%%%%%%%%%%%%%%%%%%%%%%%%%%%%%%%%%%%%%%%%%
\subsection{Valid Inequalities}
\label{subsec:valid-inequalities}
%%%%%%%%%%%%%%%%%%%%%%%%%%%%%%%%%%%%%%%%%%%%%%%%%%%%%%%%%%%%

An inequality is valid for an integer feasible set \(S\) if

\[
\boxed{
a^Tx\leq b
}
\]

for every

\[
x\in S.
\]

Adding valid inequalities does not change the integer feasible solutions.

It only strengthens the continuous relaxation.

%%%%%%%%%%%%%%%%%%%%%%%%%%%%%%%%%%%%%%%%%%%%%%%%%%%%%%%%%%%%
\subsection{Geometric Meaning of a Cut}
\label{subsec:geometric-cutting-plane}
%%%%%%%%%%%%%%%%%%%%%%%%%%%%%%%%%%%%%%%%%%%%%%%%%%%%%%%%%%%%

The LP relaxation may contain fractional vertices outside the integer
hull.

A cutting plane removes part of the relaxation while preserving

\[
\boxed{
\operatorname{conv}
(
P\cap\mathbb{Z}^n
).
}
\]

Ideally, repeated cuts make the relaxation increasingly resemble the
integer hull.

%%%%%%%%%%%%%%%%%%%%%%%%%%%%%%%%%%%%%%%%%%%%%%%%%%%%%%%%%%%%
\subsection{Gomory Cuts Preview}
\label{subsec:gomory-cuts}
%%%%%%%%%%%%%%%%%%%%%%%%%%%%%%%%%%%%%%%%%%%%%%%%%%%%%%%%%%%%

Gomory cutting planes can be derived from fractional rows of a simplex
tableau.

The basic idea is to exploit the fact that certain variables must be
integer.

Fractional relationships imply additional valid inequalities that exclude
the current fractional basic solution.

%%%%%%%%%%%%%%%%%%%%%%%%%%%%%%%%%%%%%%%%%%%%%%%%%%%%%%%%%%%%
\subsection{Cover Inequalities}
\label{subsec:cover-inequalities}
%%%%%%%%%%%%%%%%%%%%%%%%%%%%%%%%%%%%%%%%%%%%%%%%%%%%%%%%%%%%

Consider a binary knapsack constraint

\[
\boxed{
\sum_{j=1}^{n}
a_jx_j
\leq
b,
}
\]

with

\[
x_j\in\{0,1\}.
\]

A set \(C\) is a cover if

\[
\boxed{
\sum_{j\in C}
a_j
>
b.
}
\]

All items in \(C\) cannot be selected simultaneously.

Therefore,

\[
\boxed{
\sum_{j\in C}
x_j
\leq
|C|-1
}
\]

is a valid inequality.

%%%%%%%%%%%%%%%%%%%%%%%%%%%%%%%%%%%%%%%%%%%%%%%%%%%%%%%%%%%%
\subsection{Worked Example: Cover Cut}
\label{subsec:worked-cover-cut}
%%%%%%%%%%%%%%%%%%%%%%%%%%%%%%%%%%%%%%%%%%%%%%%%%%%%%%%%%%%%

\begin{workedexamplebox}{Binary Knapsack Cover}

Consider

\[
6x_1+5x_2+4x_3\leq10,
\]

with

\[
x_j\in\{0,1\}.
\]

The set

\[
C=\{1,2\}
\]

is a cover because

\[
6+5=11>10.
\]

Therefore \(x_1\) and \(x_2\) cannot both equal \(1\).

A valid cover inequality is

\begin{answerbox}
\[
\boxed{
x_1+x_2\leq1.
}
\]
\end{answerbox}

\end{workedexamplebox}

%%%%%%%%%%%%%%%%%%%%%%%%%%%%%%%%%%%%%%%%%%%%%%%%%%%%%%%%%%%%
\subsection{Branch-and-Cut}
\label{subsec:branch-and-cut}
%%%%%%%%%%%%%%%%%%%%%%%%%%%%%%%%%%%%%%%%%%%%%%%%%%%%%%%%%%%%

Modern MILP solvers combine branch-and-bound with cutting planes.

The resulting framework is called branch-and-cut.

At nodes, the solver may:

\[
\boxed{
\text{solve an LP relaxation},
}
\]

\[
\boxed{
\text{generate cuts},
}
\]

\[
\boxed{
\text{search for integer solutions},
}
\]

and

\[
\boxed{
\text{branch if needed}.
}
\]

%%%%%%%%%%%%%%%%%%%%%%%%%%%%%%%%%%%%%%%%%%%%%%%%%%%%%%%%%%%%
\subsection{Cut Separation}
\label{subsec:cut-separation}
%%%%%%%%%%%%%%%%%%%%%%%%%%%%%%%%%%%%%%%%%%%%%%%%%%%%%%%%%%%%

A cut-separation problem asks:

\[
\boxed{
\text{Does the current fractional solution violate a useful valid inequality?}
}
\]

If so, generate one or more such inequalities.

Efficient separation algorithms are important for many structured
integer programs.

%%%%%%%%%%%%%%%%%%%%%%%%%%%%%%%%%%%%%%%%%%%%%%%%%%%%%%%%%%%%
\subsection{Primal Heuristics}
\label{subsec:integer-primal-heuristics}
%%%%%%%%%%%%%%%%%%%%%%%%%%%%%%%%%%%%%%%%%%%%%%%%%%%%%%%%%%%%

A primal heuristic attempts to find good integer feasible solutions
quickly.

Examples include:

\[
\boxed{
\text{rounding heuristics},
}
\]

\[
\boxed{
\text{repair heuristics},
}
\]

\[
\boxed{
\text{local search},
}
\]

and

\[
\boxed{
\text{relaxation-guided heuristics}.
}
\]

A heuristic solution provides an incumbent but not, by itself, a proof of
optimality.

%%%%%%%%%%%%%%%%%%%%%%%%%%%%%%%%%%%%%%%%%%%%%%%%%%%%%%%%%%%%
\subsection{Feasibility Pump Preview}
\label{subsec:feasibility-pump-preview}
%%%%%%%%%%%%%%%%%%%%%%%%%%%%%%%%%%%%%%%%%%%%%%%%%%%%%%%%%%%%

The feasibility pump alternates between:

\[
\boxed{
\text{solutions of an LP-type relaxation}
}
\]

and

\[
\boxed{
\text{nearby integer points}.
}
\]

Its goal is to find an integer feasible point rapidly.

It is a heuristic and therefore does not independently certify global
optimality.

%%%%%%%%%%%%%%%%%%%%%%%%%%%%%%%%%%%%%%%%%%%%%%%%%%%%%%%%%%%%
\subsection{Local Branching Preview}
\label{subsec:local-branching-preview}
%%%%%%%%%%%%%%%%%%%%%%%%%%%%%%%%%%%%%%%%%%%%%%%%%%%%%%%%%%%%

Given an incumbent binary solution \(\bar{x}\), local branching restricts
search to solutions differing in only a limited number of binary
components.

A Hamming-distance neighborhood can be imposed using a constraint of the
form

\[
\boxed{
\sum_{j:\bar{x}_j=0}
x_j
+
\sum_{j:\bar{x}_j=1}
(1-x_j)
\leq
k.
}
\]

This defines a neighborhood around the incumbent.

%%%%%%%%%%%%%%%%%%%%%%%%%%%%%%%%%%%%%%%%%%%%%%%%%%%%%%%%%%%%
\subsection{Presolve}
\label{subsec:integer-presolve}
%%%%%%%%%%%%%%%%%%%%%%%%%%%%%%%%%%%%%%%%%%%%%%%%%%%%%%%%%%%%

Presolve simplifies the model before and during optimization.

Possible reductions include:

\[
\boxed{
\text{fixing variables},
}
\]

\[
\boxed{
\text{removing redundant constraints},
}
\]

\[
\boxed{
\text{tightening bounds},
}
\]

\[
\boxed{
\text{detecting implied relations},
}
\]

and

\[
\boxed{
\text{detecting infeasibility}.
}
\]

Presolve can dramatically reduce search-tree size.

%%%%%%%%%%%%%%%%%%%%%%%%%%%%%%%%%%%%%%%%%%%%%%%%%%%%%%%%%%%%
\subsection{Bound Tightening}
\label{subsec:integer-bound-tightening}
%%%%%%%%%%%%%%%%%%%%%%%%%%%%%%%%%%%%%%%%%%%%%%%%%%%%%%%%%%%%

Suppose

\[
x\leq100
\]

is initially known.

Other constraints may imply the stronger bound

\[
x\leq17.
\]

Replacing \(100\) by \(17\) can improve the LP relaxation and any Big-\(M\)
constraints involving \(x\).

Tight bounds are highly valuable in mixed-integer modeling.

%%%%%%%%%%%%%%%%%%%%%%%%%%%%%%%%%%%%%%%%%%%%%%%%%%%%%%%%%%%%
\subsection{Implied Bounds}
\label{subsec:implied-bounds}
%%%%%%%%%%%%%%%%%%%%%%%%%%%%%%%%%%%%%%%%%%%%%%%%%%%%%%%%%%%%

An implied bound follows from other constraints.

For example, if

\[
2x+3y\leq12,
\]

and

\[
y\geq0,
\]

then

\[
\boxed{
x\leq6.
}
\]

If \(x\) must be integer,

\[
\boxed{
x\leq6
}
\]

remains valid and may strengthen later formulations.

%%%%%%%%%%%%%%%%%%%%%%%%%%%%%%%%%%%%%%%%%%%%%%%%%%%%%%%%%%%%
\subsection{Symmetry}
\label{subsec:integer-symmetry}
%%%%%%%%%%%%%%%%%%%%%%%%%%%%%%%%%%%%%%%%%%%%%%%%%%%%%%%%%%%%

An integer program has symmetry when several assignments of variables
represent mathematically equivalent solutions.

Symmetry can cause branch-and-bound to explore many redundant portions of
the search tree.

%%%%%%%%%%%%%%%%%%%%%%%%%%%%%%%%%%%%%%%%%%%%%%%%%%%%%%%%%%%%
\subsection{Symmetry-Breaking Constraints}
\label{subsec:symmetry-breaking}
%%%%%%%%%%%%%%%%%%%%%%%%%%%%%%%%%%%%%%%%%%%%%%%%%%%%%%%%%%%%

Suppose identical facilities have binary variables

\[
y_1,y_2,\ldots,y_m.
\]

One may impose

\[
\boxed{
y_1\geq y_2\geq\cdots\geq y_m
}
\]

when this ordering does not remove genuinely distinct solutions.

This can eliminate equivalent permutations.

%%%%%%%%%%%%%%%%%%%%%%%%%%%%%%%%%%%%%%%%%%%%%%%%%%%%%%%%%%%%
\subsection{Total Unimodularity}
\label{subsec:integer-total-unimodularity}
%%%%%%%%%%%%%%%%%%%%%%%%%%%%%%%%%%%%%%%%%%%%%%%%%%%%%%%%%%%%

Recall that a matrix is totally unimodular if every square subdeterminant
belongs to

\[
\boxed{
\{-1,0,1\}.
}
\]

For important linear systems, total unimodularity combined with integral
right-hand-side data guarantees integral extreme points.

Thus an integer restriction may become unnecessary.

%%%%%%%%%%%%%%%%%%%%%%%%%%%%%%%%%%%%%%%%%%%%%%%%%%%%%%%%%%%%
\subsection{Automatic Integrality}
\label{subsec:automatic-integrality}
%%%%%%%%%%%%%%%%%%%%%%%%%%%%%%%%%%%%%%%%%%%%%%%%%%%%%%%%%%%%

Suppose

\[
A
\]

is totally unimodular and

\[
b
\]

is integral.

Under suitable standard-form assumptions, vertices of

\[
\boxed{
\{
x:
Ax\leq b
\}
}
\]

are integral.

This explains why many:

\[
\boxed{
\text{network-flow},
}
\]

\[
\boxed{
\text{assignment},
}
\]

and

\[
\boxed{
\text{transportation}
}
\]

problems can be solved efficiently as linear programs.

%%%%%%%%%%%%%%%%%%%%%%%%%%%%%%%%%%%%%%%%%%%%%%%%%%%%%%%%%%%%
\subsection{Extended Formulations}
\label{subsec:integer-extended-formulations}
%%%%%%%%%%%%%%%%%%%%%%%%%%%%%%%%%%%%%%%%%%%%%%%%%%%%%%%%%%%%

Sometimes adding variables and constraints produces a higher-dimensional
formulation with a much stronger LP relaxation.

This is called an extended formulation.

Although it contains more variables, it may greatly reduce
branch-and-bound effort.

Thus:

\[
\boxed{
\text{larger formulation}
}
\]

can sometimes mean

\[
\boxed{
\text{easier optimization}.
}
\]

%%%%%%%%%%%%%%%%%%%%%%%%%%%%%%%%%%%%%%%%%%%%%%%%%%%%%%%%%%%%
\subsection{Formulation Strength}
\label{subsec:integer-formulation-strength}
%%%%%%%%%%%%%%%%%%%%%%%%%%%%%%%%%%%%%%%%%%%%%%%%%%%%%%%%%%%%

Two MILP formulations may represent exactly the same integer feasible
solutions but have different LP relaxations.

A stronger formulation generally excludes more fractional points.

Therefore formulation comparison should examine:

\[
\boxed{
\text{relaxation bounds},
}
\]

not merely the number of constraints.

%%%%%%%%%%%%%%%%%%%%%%%%%%%%%%%%%%%%%%%%%%%%%%%%%%%%%%%%%%%%
\subsection{Lagrangian Relaxation}
\label{subsec:lagrangian-relaxation-integer}
%%%%%%%%%%%%%%%%%%%%%%%%%%%%%%%%%%%%%%%%%%%%%%%%%%%%%%%%%%%%

Suppose certain constraints make an integer problem difficult.

For minimization,

\[
Ax=b
\]

may be moved into the objective using multipliers \(\lambda\):

\[
\boxed{
L(x,\lambda)
=
c^Tx
+
\lambda^T(Ax-b).
}
\]

The resulting relaxed problem may decompose into easier subproblems.

Optimizing over the multipliers can produce useful bounds.

%%%%%%%%%%%%%%%%%%%%%%%%%%%%%%%%%%%%%%%%%%%%%%%%%%%%%%%%%%%%
\subsection{Lagrangian Bounds}
\label{subsec:lagrangian-bounds-integer}
%%%%%%%%%%%%%%%%%%%%%%%%%%%%%%%%%%%%%%%%%%%%%%%%%%%%%%%%%%%%

For a minimization problem, a properly constructed Lagrangian relaxation
typically provides a lower bound on the original optimum.

The multipliers are then adjusted to maximize this lower bound.

This produces a nonsmooth dual optimization problem.

%%%%%%%%%%%%%%%%%%%%%%%%%%%%%%%%%%%%%%%%%%%%%%%%%%%%%%%%%%%%
\subsection{Decomposition Preview}
\label{subsec:integer-decomposition-preview}
%%%%%%%%%%%%%%%%%%%%%%%%%%%%%%%%%%%%%%%%%%%%%%%%%%%%%%%%%%%%

Large integer programs may contain separable structure.

Decomposition techniques attempt to exploit this by dividing the problem
into:

\[
\boxed{
\text{master decisions}
}
\]

and

\[
\boxed{
\text{subproblems}.
}
\]

Important approaches include:

\[
\boxed{
\text{Dantzig--Wolfe decomposition},
}
\]

\[
\boxed{
\text{column generation},
}
\]

and

\[
\boxed{
\text{Benders decomposition}.
}
\]

%%%%%%%%%%%%%%%%%%%%%%%%%%%%%%%%%%%%%%%%%%%%%%%%%%%%%%%%%%%%
\subsection{Benders Decomposition Preview}
\label{subsec:benders-decomposition-preview}
%%%%%%%%%%%%%%%%%%%%%%%%%%%%%%%%%%%%%%%%%%%%%%%%%%%%%%%%%%%%

Suppose binary design variables \(y\) interact with continuous operating
variables \(x\).

A Benders approach may separate:

\[
\boxed{
\text{master problem in }y
}
\]

from

\[
\boxed{
\text{continuous subproblem in }x.
}
\]

Information from the subproblem is returned to the master through Benders
cuts.

This structure is especially useful in facility and network design.

%%%%%%%%%%%%%%%%%%%%%%%%%%%%%%%%%%%%%%%%%%%%%%%%%%%%%%%%%%%%
\subsection{Column Generation Preview}
\label{subsec:integer-column-generation-preview}
%%%%%%%%%%%%%%%%%%%%%%%%%%%%%%%%%%%%%%%%%%%%%%%%%%%%%%%%%%%%

Some formulations contain an enormous number of possible variables.

Column generation begins with only a small subset.

A pricing subproblem searches for a new variable with attractive reduced
cost.

This idea becomes the foundation of branch-and-price when integrality is
also required.

%%%%%%%%%%%%%%%%%%%%%%%%%%%%%%%%%%%%%%%%%%%%%%%%%%%%%%%%%%%%
\subsection{Branch-and-Price Preview}
\label{subsec:branch-and-price-preview}
%%%%%%%%%%%%%%%%%%%%%%%%%%%%%%%%%%%%%%%%%%%%%%%%%%%%%%%%%%%%

Branch-and-price combines:

\[
\boxed{
\text{branch-and-bound}
}
\]

with

\[
\boxed{
\text{column generation}.
}
\]

It is useful when the natural LP relaxation has too many columns to list
explicitly.

Applications include:

\[
\boxed{
\text{vehicle routing},
}
\]

\[
\boxed{
\text{crew scheduling},
}
\]

and

\[
\boxed{
\text{cutting-stock problems}.
}
\]

%%%%%%%%%%%%%%%%%%%%%%%%%%%%%%%%%%%%%%%%%%%%%%%%%%%%%%%%%%%%
\subsection{Complexity of Integer Programming}
\label{subsec:integer-programming-complexity}
%%%%%%%%%%%%%%%%%%%%%%%%%%%%%%%%%%%%%%%%%%%%%%%%%%%%%%%%%%%%

General integer programming is computationally much harder than linear
programming.

Many important integer-programming problems are NP-hard.

Therefore large instances may require:

\[
\boxed{
\text{strong formulations},
}
\]

\[
\boxed{
\text{good bounds},
}
\]

\[
\boxed{
\text{problem-specific cuts},
}
\]

and

\[
\boxed{
\text{effective heuristics}.
}
\]

%%%%%%%%%%%%%%%%%%%%%%%%%%%%%%%%%%%%%%%%%%%%%%%%%%%%%%%%%%%%
\subsection{Pseudo-Polynomial Algorithms}
\label{subsec:pseudo-polynomial-integer}
%%%%%%%%%%%%%%%%%%%%%%%%%%%%%%%%%%%%%%%%%%%%%%%%%%%%%%%%%%%%

Some integer problems admit dynamic-programming algorithms whose running
time depends on the numerical size of the data rather than only on the
number of input bits.

The classical knapsack problem has such algorithms.

These are called pseudo-polynomial algorithms.

This distinction becomes important in computational complexity.

%%%%%%%%%%%%%%%%%%%%%%%%%%%%%%%%%%%%%%%%%%%%%%%%%%%%%%%%%%%%
\subsection{Primal and Dual Information}
\label{subsec:mip-primal-dual-information}
%%%%%%%%%%%%%%%%%%%%%%%%%%%%%%%%%%%%%%%%%%%%%%%%%%%%%%%%%%%%

In MILP terminology:

\[
\boxed{
\text{primal solution}
=
\text{integer feasible solution},
}
\]

while a relaxation gives a bound analogous to dual information.

For minimization,

\[
\boxed{
\text{incumbent}
=
\text{upper bound},
}
\]

and

\[
\boxed{
\text{best node bound}
=
\text{lower bound}.
}
\]

Optimization ends when these bounds are sufficiently close.

%%%%%%%%%%%%%%%%%%%%%%%%%%%%%%%%%%%%%%%%%%%%%%%%%%%%%%%%%%%%
\subsection{Optimality Certificate}
\label{subsec:mip-optimality-certificate}
%%%%%%%%%%%%%%%%%%%%%%%%%%%%%%%%%%%%%%%%%%%%%%%%%%%%%%%%%%%%

Finding a feasible integer solution does not prove optimality.

An optimality certificate requires a valid global bound.

For minimization, if

\[
\boxed{
U=L,
}
\]

where \(U\) is the incumbent value and \(L\) is the global lower bound,
then

\[
\boxed{
z_{\mathrm{IP}}^*=U=L.
}
\]

%%%%%%%%%%%%%%%%%%%%%%%%%%%%%%%%%%%%%%%%%%%%%%%%%%%%%%%%%%%%
\subsection{Time Limits and MIP Gaps}
\label{subsec:mip-time-limit-gap}
%%%%%%%%%%%%%%%%%%%%%%%%%%%%%%%%%%%%%%%%%%%%%%%%%%%%%%%%%%%%

Large MILPs may be stopped before exact optimality is proven.

Suppose a minimization solver stops with

\[
U=102,
\]

and

\[
L=100.
\]

Then the true optimum satisfies

\[
\boxed{
100
\leq
z^*
\leq
102.
}
\]

Even without an exact proof, the bound quantifies solution quality.

%%%%%%%%%%%%%%%%%%%%%%%%%%%%%%%%%%%%%%%%%%%%%%%%%%%%%%%%%%%%
\subsection{Worked Example: Bound Interpretation}
\label{subsec:worked-mip-bound-interpretation}
%%%%%%%%%%%%%%%%%%%%%%%%%%%%%%%%%%%%%%%%%%%%%%%%%%%%%%%%%%%%

\begin{workedexamplebox}{Reading Solver Bounds}

Suppose a minimization MILP has:

\[
\text{incumbent}=500,
\]

and

\[
\text{best bound}=490.
\]

Then

\[
490
\leq
z^*
\leq
500.
\]

The absolute gap is

\[
500-490=10.
\]

Therefore,

\begin{answerbox}
\[
\boxed{
490
\leq
z^*
\leq
500.
}
\]
\end{answerbox}

The incumbent is feasible, while the lower bound certifies that no
solution below \(490\) can exist.

\end{workedexamplebox}

%%%%%%%%%%%%%%%%%%%%%%%%%%%%%%%%%%%%%%%%%%%%%%%%%%%%%%%%%%%%
\subsection{Modeling Workflow}
\label{subsec:integer-programming-workflow}
%%%%%%%%%%%%%%%%%%%%%%%%%%%%%%%%%%%%%%%%%%%%%%%%%%%%%%%%%%%%

A practical integer-programming workflow is:

\begin{enumerate}
    \item identify which decisions are inherently discrete;
    \item define integer and binary variables clearly;
    \item formulate the objective;
    \item write physical and resource constraints;
    \item encode logical relations;
    \item derive tight variable bounds;
    \item avoid unnecessary Big-\(M\) constants;
    \item inspect the LP relaxation;
    \item strengthen the formulation when possible;
    \item exploit known structure such as networks or assignments;
    \item solve with branch-and-bound or branch-and-cut;
    \item inspect the incumbent and global bound;
    \item inspect the optimality gap;
    \item analyze sensitivity carefully;
    \item interpret the integer solution in application terms.
\end{enumerate}

%%%%%%%%%%%%%%%%%%%%%%%%%%%%%%%%%%%%%%%%%%%%%%%%%%%%%%%%%%%%
\subsection{Choosing an Integer-Programming Approach}
\label{subsec:choosing-integer-programming-method}
%%%%%%%%%%%%%%%%%%%%%%%%%%%%%%%%%%%%%%%%%%%%%%%%%%%%%%%%%%%%

For general MILPs, use:

\[
\boxed{
\text{branch-and-cut}.
}
\]

For strong network structure, consider:

\[
\boxed{
\text{specialized network algorithms}.
}
\]

For enormous variable sets, consider:

\[
\boxed{
\text{column generation or branch-and-price}.
}
\]

For complicating linking constraints, consider:

\[
\boxed{
\text{Lagrangian or Benders decomposition}.
}
\]

For small structured problems, dynamic programming may be effective.

%%%%%%%%%%%%%%%%%%%%%%%%%%%%%%%%%%%%%%%%%%%%%%%%%%%%%%%%%%%%
\subsection{Common Errors}
\label{subsec:integer-programming-common-errors}
%%%%%%%%%%%%%%%%%%%%%%%%%%%%%%%%%%%%%%%%%%%%%%%%%%%%%%%%%%%%

\subsubsection{Using Integer Variables When They Are Not Necessary}

If a model's LP structure already guarantees integral extreme points,
explicit integer restrictions can create unnecessary computational cost.

Check for network or total-unimodular structure.

%%%%%%%%%%%%%%%%%%%%%%%%%%%%%%%%%%%%%%%%%%%%%%%%%%%%%%%%%%%%
\subsubsection{Treating Binary Variables as Continuous Probabilities}

A binary decision variable represents

\[
\boxed{
0
\text{ or }
1,
}
\]

not an arbitrary fraction between them.

Fractional values arise only in relaxations.

%%%%%%%%%%%%%%%%%%%%%%%%%%%%%%%%%%%%%%%%%%%%%%%%%%%%%%%%%%%%
\subsubsection{Rounding Without Checking Feasibility}

An LP solution cannot generally be rounded componentwise.

The rounded point may violate constraints or destroy solution quality.

%%%%%%%%%%%%%%%%%%%%%%%%%%%%%%%%%%%%%%%%%%%%%%%%%%%%%%%%%%%%
\subsubsection{Using an Invalid Big-\(M\)}

If \(M\) is too small, valid feasible solutions may be incorrectly
removed.

The formulation is then mathematically wrong.

%%%%%%%%%%%%%%%%%%%%%%%%%%%%%%%%%%%%%%%%%%%%%%%%%%%%%%%%%%%%
\subsubsection{Using an Excessively Large Big-\(M\)}

If \(M\) is far larger than necessary, the LP relaxation becomes weak and
numerical conditioning can deteriorate.

Use tight valid bounds.

%%%%%%%%%%%%%%%%%%%%%%%%%%%%%%%%%%%%%%%%%%%%%%%%%%%%%%%%%%%%
\subsubsection{Forgetting to Link Binary and Continuous Variables}

A fixed-charge variable \(y\) has no effect unless constraints such as

\[
\boxed{
x\leq Uy
}
\]

connect it to the activity quantity.

%%%%%%%%%%%%%%%%%%%%%%%%%%%%%%%%%%%%%%%%%%%%%%%%%%%%%%%%%%%%
\subsubsection{Confusing Feasible Integer Solutions with Optimal Solutions}

A heuristic may find a good feasible solution quickly.

Without a matching global bound, it does not prove optimality.

%%%%%%%%%%%%%%%%%%%%%%%%%%%%%%%%%%%%%%%%%%%%%%%%%%%%%%%%%%%%
\subsubsection{Confusing LP Bounds in Minimization and Maximization}

For minimization,

\[
\boxed{
z_{\mathrm{LP}}^*
\leq
z_{\mathrm{IP}}^*.
}
\]

For maximization,

\[
\boxed{
z_{\mathrm{IP}}^*
\leq
z_{\mathrm{LP}}^*.
}
\]

The direction of the bound reverses.

%%%%%%%%%%%%%%%%%%%%%%%%%%%%%%%%%%%%%%%%%%%%%%%%%%%%%%%%%%%%
\subsubsection{Branching Without Covering All Integer Possibilities}

If

\[
x_j^*=3.7,
\]

the correct branches are

\[
x_j\leq3
\]

and

\[
x_j\geq4.
\]

Together they include every integer possibility.

%%%%%%%%%%%%%%%%%%%%%%%%%%%%%%%%%%%%%%%%%%%%%%%%%%%%%%%%%%%%
\subsubsection{Pruning a Node Without a Valid Bound}

A node may only be discarded by a logically valid reason such as:

\[
\boxed{
\text{infeasibility},
}
\]

\[
\boxed{
\text{integrality},
}
\]

or

\[
\boxed{
\text{a bound proving it cannot improve the incumbent}.
}
\]

%%%%%%%%%%%%%%%%%%%%%%%%%%%%%%%%%%%%%%%%%%%%%%%%%%%%%%%%%%%%
\subsubsection{Assuming More Constraints Always Make a Model Worse}

Useful valid inequalities can make the model easier by strengthening the
LP relaxation.

Constraint count alone is not a measure of formulation quality.

%%%%%%%%%%%%%%%%%%%%%%%%%%%%%%%%%%%%%%%%%%%%%%%%%%%%%%%%%%%%
\subsubsection{Ignoring Symmetry}

Equivalent binary configurations can create enormous redundant search
trees.

Symmetry-breaking constraints may dramatically improve performance.

%%%%%%%%%%%%%%%%%%%%%%%%%%%%%%%%%%%%%%%%%%%%%%%%%%%%%%%%%%%%
\subsubsection{Treating Solver Gap as Constraint Violation}

The MIP optimality gap compares the incumbent objective with a global
bound.

It is not the same thing as primal constraint violation.

%%%%%%%%%%%%%%%%%%%%%%%%%%%%%%%%%%%%%%%%%%%%%%%%%%%%%%%%%%%%
\subsubsection{Interpreting a Time-Limited Incumbent as Proven Optimal}

A time-limited solver may return an excellent feasible solution without
closing the global bound.

Always inspect the reported gap and termination status.

%%%%%%%%%%%%%%%%%%%%%%%%%%%%%%%%%%%%%%%%%%%%%%%%%%%%%%%%%%%%
\subsection{Applications}
\label{subsec:integer-programming-applications}
%%%%%%%%%%%%%%%%%%%%%%%%%%%%%%%%%%%%%%%%%%%%%%%%%%%%%%%%%%%%

\begin{applicationbox}

\textbf{Facility Location.}
Binary variables determine which warehouses, plants, hospitals, or
service centers should be opened.

\medskip

\textbf{Scheduling.}
Integer decisions assign jobs to machines, workers to shifts, and tasks
to time periods.

\medskip

\textbf{Transportation and Logistics.}
Mixed-integer models determine routes, fleets, shipment choices, and
distribution-network designs.

\medskip

\textbf{Production Planning.}
Setup decisions, minimum batch sizes, capacity expansion, and product
selection naturally require binary or integer variables.

\medskip

\textbf{Finance.}
Portfolio models may include minimum transaction sizes, cardinality
limits, or discrete asset-selection decisions.

\medskip

\textbf{Energy Systems.}
Unit commitment determines which generators are switched on or off over
time.

\medskip

\textbf{Telecommunications.}
Network design chooses discrete links, capacities, and equipment
placements.

\medskip

\textbf{Public Planning.}
Emergency facilities, staffing levels, school locations, and resource
deployment can be modeled using MILPs.

\end{applicationbox}

%%%%%%%%%%%%%%%%%%%%%%%%%%%%%%%%%%%%%%%%%%%%%%%%%%%%%%%%%%%%
\subsection{Exercises}
\label{subsec:integer-programming-exercises}
%%%%%%%%%%%%%%%%%%%%%%%%%%%%%%%%%%%%%%%%%%%%%%%%%%%%%%%%%%%%

\begin{exercise}[1]
Define an integer program.
\end{exercise}

\begin{exercise}[1]
Define a binary variable.
\end{exercise}

\begin{exercise}[1]
Define a mixed-integer linear program.
\end{exercise}

\begin{exercise}[1]
Define an LP relaxation.
\end{exercise}

\begin{exercise}[1]
Define the integer hull.
\end{exercise}

\begin{exercise}[1]
Define an incumbent.
\end{exercise}

\begin{exercise}[1]
Define branch-and-bound.
\end{exercise}

\begin{exercise}[1]
Define a valid inequality.
\end{exercise}

\begin{exercise}[1]
Define a cutting plane.
\end{exercise}

\begin{exercise}[1]
Define the integrality gap conceptually.
\end{exercise}

\begin{exercise}[2]
Write a binary constraint requiring exactly one of

\[
x_1,x_2,x_3
\]

to equal \(1\).
\end{exercise}

\begin{exercise}[2]
Write a constraint requiring at most two of five binary variables to equal
\(1\).
\end{exercise}

\begin{exercise}[2]
Model

\[
x=1
\Rightarrow
y=1
\]

using binary variables.
\end{exercise}

\begin{exercise}[2]
Model mutual exclusion between binary variables \(x\) and \(y\).
\end{exercise}

\begin{exercise}[2]
Formulate a fixed-charge relation between continuous variable \(q\) and
binary variable \(y\) when

\[
0\leq q\leq100
\]

whenever the activity is enabled.
\end{exercise}

\begin{exercise}[2]
Write the LP relaxation of

\[
\begin{aligned}
\max\quad&
5x_1+4x_2
\\
\text{subject to}\quad&
3x_1+2x_2\leq9,
\\
&
x_1,x_2\in\mathbb{Z}_{\geq0}.
\end{aligned}
\]
\end{exercise}

\begin{exercise}[2]
Explain why the LP relaxation gives an upper bound for the previous
maximization problem.
\end{exercise}

\begin{exercise}[2]
Suppose an LP relaxation gives

\[
x_1=4.3.
\]

Write the two standard branch constraints.
\end{exercise}

\begin{exercise}[2]
Explain three reasons a branch-and-bound node may be pruned.
\end{exercise}

\begin{exercise}[3]
Formulate a binary knapsack model with four items.
\end{exercise}

\begin{exercise}[3]
Solve a small binary knapsack problem by complete enumeration.
\end{exercise}

\begin{exercise}[3]
Compare the integer optimum with the LP-relaxation optimum for the
previous problem.
\end{exercise}

\begin{exercise}[3]
Construct an assignment model for three workers and three jobs.
\end{exercise}

\begin{exercise}[3]
Formulate a facility-location problem with two possible facilities and
three customers.
\end{exercise}

\begin{exercise}[3]
Add capacity constraints to the previous model.
\end{exercise}

\begin{exercise}[3]
Formulate a setup-cost production model.
\end{exercise}

\begin{exercise}[3]
Model the implication

\[
y=1
\Rightarrow
2x_1+3x_2\leq10
\]

using Big-\(M\).
\end{exercise}

\begin{exercise}[3]
Explain how a valid upper bound on \(x_1,x_2\) can be used to derive a
tight \(M\).
\end{exercise}

\begin{exercise}[3]
Formulate two nonoverlapping jobs using one binary sequencing variable.
\end{exercise}

\begin{exercise}[3]
For a minimization problem with incumbent value \(80\) and node lower
bound \(85\), determine whether the node can be pruned.
\end{exercise}

\begin{exercise}[3]
For a minimization problem with incumbent \(120\) and global lower bound
\(114\), compute the absolute gap.
\end{exercise}

\begin{exercise}[3]
Explain why a fractional LP solution can still provide valuable
information.
\end{exercise}

\begin{exercise}[3]
Find a cover inequality for

\[
7x_1+6x_2+4x_3\leq10,
\]

with binary variables.
\end{exercise}

\begin{exercise}[3]
Explain how a cutting plane can improve branch-and-bound.
\end{exercise}

\begin{exercise}[3]
Explain the difference between a primal heuristic and a cutting plane.
\end{exercise}

\begin{exercise}[3]
Give an example of symmetry in a binary model.
\end{exercise}

\begin{exercise}[3]
Construct a symmetry-breaking constraint for three interchangeable
facilities.
\end{exercise}

\begin{exercise}[3]
Explain how bound tightening improves a Big-\(M\) formulation.
\end{exercise}

\begin{exercise}[3]
Explain why the assignment problem can often be solved using its LP
relaxation.
\end{exercise}

\begin{exercise}[3]
Interpret a solver report with incumbent \(250\) and best lower bound
\(247\) for a minimization problem.
\end{exercise}

\begin{exercise}[3]
Explain why stopping with a small MIP gap may be acceptable in a large
practical model.
\end{exercise}

\begin{exercise}[4]
Study the integer hull of a two-dimensional integer program.
\end{exercise}

\begin{exercise}[4]
Derive Gomory fractional cuts from a simplex tableau.
\end{exercise}

\begin{exercise}[4]
Study mixed-integer rounding cuts.
\end{exercise}

\begin{exercise}[4]
Study lifted cover inequalities.
\end{exercise}

\begin{exercise}[4]
Investigate branch-and-cut algorithms.
\end{exercise}

\begin{exercise}[4]
Compare branching-variable selection strategies.
\end{exercise}

\begin{exercise}[4]
Study strong branching and pseudocost branching.
\end{exercise}

\begin{exercise}[4]
Investigate node-selection strategies in branch-and-bound.
\end{exercise}

\begin{exercise}[4]
Study total unimodularity and integral polyhedra.
\end{exercise}

\begin{exercise}[4]
Prove integrality of the assignment polytope.
\end{exercise}

\begin{exercise}[4]
Study Lagrangian relaxation for integer programming.
\end{exercise}

\begin{exercise}[4]
Derive a Lagrangian dual bound for a knapsack-type model.
\end{exercise}

\begin{exercise}[4]
Study Benders decomposition for facility-location models.
\end{exercise}

\begin{exercise}[4]
Study Dantzig--Wolfe decomposition.
\end{exercise}

\begin{exercise}[4]
Study column generation for cutting-stock problems.
\end{exercise}

\begin{exercise}[4]
Investigate branch-and-price.
\end{exercise}

\begin{exercise}[4]
Study symmetry detection and orbital branching.
\end{exercise}

\begin{exercise}[4]
Compare weak and strong MILP formulations for the same logical model.
\end{exercise}

\begin{exercise}[4]
Study extended formulations and polyhedral complexity.
\end{exercise}

\begin{exercise}[4]
Investigate computational complexity of general integer programming.
\end{exercise}

%%%%%%%%%%%%%%%%%%%%%%%%%%%%%%%%%%%%%%%%%%%%%%%%%%%%%%%%%%%%
\subsection{Conceptual Summary}
\label{subsec:integer-programming-summary}
%%%%%%%%%%%%%%%%%%%%%%%%%%%%%%%%%%%%%%%%%%%%%%%%%%%%%%%%%%%%

An integer program imposes restrictions such as

\[
\boxed{
x_j\in\mathbb{Z},
}
\]

or

\[
\boxed{
x_j\in\{0,1\}.
}
\]

A mixed-integer linear program combines integer and continuous variables:

\[
\boxed{
\begin{aligned}
\min
\quad&
c^Tx+d^Ty
\\
\text{subject to}
\quad&
Ax+By\leq b,
\\
&
x\in\mathbb{Z}^p,
\\
&
y\in\mathbb{R}^q.
\end{aligned}
}
\]

Removing integrality gives the LP relaxation.

For minimization,

\[
\boxed{
z_{\mathrm{LP}}^*
\leq
z_{\mathrm{IP}}^*.
}
\]

For maximization,

\[
\boxed{
z_{\mathrm{IP}}^*
\leq
z_{\mathrm{LP}}^*.
}
\]

Branch-and-bound repeatedly:

\[
\boxed{
\text{relaxes},
}
\]

\[
\boxed{
\text{branches},
}
\]

and

\[
\boxed{
\text{prunes}.
}
\]

Cutting planes strengthen relaxations by removing fractional points while
preserving all integer feasible points.

Modern MILP methods combine these ideas through

\[
\boxed{
\text{branch-and-cut}.
}
\]

Binary variables allow optimization models to represent:

\[
\boxed{
\text{selection}
+
\text{logic}
+
\text{activation}
+
\text{sequencing}.
}
\]

Integer programming therefore combines

\[
\boxed{
\text{linear programming}
+
\text{discrete mathematics}
+
\text{polyhedral geometry}
+
\text{algorithms}
+
\text{decision modeling}.
}
\]

%%%%%%%%%%%%%%%%%%%%%%%%%%%%%%%%%%%%%%%%%%%%%%%%%%%%%%%%%%%%
\subsection{Section Connections}
\label{subsec:integer-programming-connections}
%%%%%%%%%%%%%%%%%%%%%%%%%%%%%%%%%%%%%%%%%%%%%%%%%%%%%%%%%%%%

\chapterconnection{
Integer Programming extends the continuous optimization framework of
Sections~\ref{sec:linear-programming} and
\ref{sec:convex-optimization} by imposing discrete decisions. LP
relaxations, integer hulls, branch-and-bound, cutting planes, and
branch-and-cut show how continuous optimization can be used to solve
discrete problems. Binary variables also provide a mathematical language
for logical and structural decisions. The next section, Combinatorial
Optimization, develops optimization directly over discrete structures
such as subsets, permutations, matchings, paths, trees, and graphs, and
connects these models with graph theory, complexity, approximation, and
specialized algorithms.
}

%%%%%%%%%%%%%%%%%%%%%%%%%%%%%%%%%%%%%%%%%%%%%%%%%%%%%%%%%%%%
% SECTION SOURCE TRACKING
%%%%%%%%%%%%%%%%%%%%%%%%%%%%%%%%%%%%%%%%%%%%%%%%%%%%%%%%%%%%

%------------------------------------------------------------
% 10.5 Integer Programming
%------------------------------------------------------------
%
% Source basis:
%   - Standard integer and mixed-integer programming theory.
%   - Standard branch-and-bound and cutting-plane methods.
%   - Standard operations-research integer modeling.
%
% Used for:
%   - pure integer programs;
%   - binary integer programs;
%   - mixed-integer linear programming;
%   - integer lattices;
%   - LP relaxations;
%   - relaxation bounds;
%   - integrality gaps;
%   - integer hulls;
%   - ideal and strong formulations;
%   - rounding limitations;
%   - binary selection and cardinality models;
%   - logical implications;
%   - mutual exclusion;
%   - either-or conditions;
%   - Big-M formulations;
%   - indicator constraints;
%   - fixed-charge models;
%   - knapsack models;
%   - assignment models;
%   - facility-location models;
%   - scheduling and sequencing formulations;
%   - branch-and-bound;
%   - branching rules;
%   - incumbents and node bounds;
%   - pruning;
%   - global and relative optimality gaps;
%   - cutting planes;
%   - valid inequalities;
%   - Gomory cuts preview;
%   - cover inequalities;
%   - branch-and-cut;
%   - cut separation;
%   - primal heuristics;
%   - feasibility pump and local branching previews;
%   - presolve;
%   - bound tightening;
%   - symmetry and symmetry breaking;
%   - total unimodularity;
%   - automatic integrality;
%   - extended formulations;
%   - Lagrangian relaxation;
%   - decomposition;
%   - Benders decomposition preview;
%   - column generation and branch-and-price previews;
%   - complexity;
%   - primal and dual bounds;
%   - solver-gap interpretation;
%   - applications and exercises.
%
% Notes:
%   - Combinatorial Optimization is developed next in Section 10.6.
%   - Dynamic Programming is developed in Section 10.7.
%   - Network Optimization is developed in Section 10.10.
%   - Scheduling Theory is developed in Section 10.12.
%
%%%%%%%%%%%%%%%%%%%%%%%%%%%%%%%%%%%%%%%%%%%%%%%%%%%%%%%%%%%%